% Moderne Titelseite
% Universität der Bundeswehr München

% === Seitennummerierung für Vorspann ===
\pagenumbering{Roman}

% === Temporäre Seitenränder für Titelseite ===
\newgeometry{
    left=2cm,
    right=2cm,
    top=2cm,
    bottom=2cm,
}

% === Titelseite ===
\begin{titlepage}
    % === Universitäts-Header ===
    \begin{center}
        \vspace*{-1.2cm}
        \IfFileExists{default_style/Signet_UniBw.pdf}{%
            \includegraphics[scale=0.95]{default_style/Signet_UniBw.pdf}%
        }{%
            \textcolor{secondary}{\rule{8cm}{3cm}}%
            \PackageWarning{titelseite}{Logo Signet_UniBw.pdf nicht gefunden}%
        }
        
        \Large
        \textsc{
            \myfaculty\\
            [0.2cm]
            \myinstitute\\
            [0.2cm]
            \myprofessur\\
            [0.2cm]
            \reviewerone
        }
    \end{center}
    
    \vfill
    
    % === Arbeitstitel ===
    \begin{center}
        \Large
        \mytype \\ % aus metadata.tex
        %[0.4cm]\Large zur Erlangung des akademischen Grades\\Bachelor of Science\\
        %[0.4cm]
        \Large
        zum Thema\\
        \vfill
        
        \huge
        \textbf{\mytitle} 
        \vfill
        
        \Large
        von\\
        [0.4cm]
        \textbf{\myname}\\
        [0.4cm]
        Matrikelnummer \matricle\\       
    \end{center}
    
    \vfill
    
    % === Betreuungs-Informationen ===
    % Leere Felder in metadata.tex (\reviewertwo, \advisor) werden ausgeblendet
    \ifdefempty{\reviewertwo}{\def\zweitgutachterzeile{}}{\def\zweitgutachterzeile{Zweitgutachter: & \reviewertwo\\}}
    \ifdefempty{\advisor}{\def\betreuerzeile{}}{\def\betreuerzeile{Betreuer: & \advisor\\}}
    \begin{center}
        \Large
        \begin{tabular}{ l l }
            Aufgabensteller: & \reviewerone\\
            \zweitgutachterzeile
            \betreuerzeile
            \\
            Tag der Einreichung: & \timeend, \myplace\\
        \end{tabular}
    \end{center}
\end{titlepage}

% === Ursprüngliche Randeinstellungen wiederherstellen ===
\restoregeometry
